\section{First implementation MVP: frozen mathematical scope}
\label{sec:mvp}

The first vNext implementation deliberately solves a smaller problem completely
before enabling every general feature in the architecture.

\subsection{MVP scene class}

The MVP supports:
\begin{enumerate}
\item $N_c\ge1$ independent superelliptic spiral coils;
\item arbitrary independent rigid poses in $\SE(3)$;
\item continuous superellipse exponent, aspect ratio, pitch, turns, conductor
      width/thickness, and superelliptic cross-section exponent;
\item finite-conductivity isotropic conductors;
\item one homogeneous, isotropic, unbounded passive electromagnetic background;
\item narrowband sinusoidal excitation at arbitrary frequency within a declared
      artifact band;
\item arbitrary complex multi-port current or voltage/circuit excitation;
\item conductor Joule loss;
\item homogeneous or piecewise-simple thermal environment sufficient for a
      Green/compact thermal transfer model;
\item linear temperature evolution first, followed by temperature-dependent
      conductor resistance as the first nonlinear extension.
\end{enumerate}

There is no restriction to coplanar coils, small rotations, matched coil shape,
or a fixed number of turns.

\subsection{Explicitly excluded from MVP}

The first implementation does not yet include:
\begin{enumerate}
\item dielectric/magnetic package scattering in the electromagnetic solve;
\item localized electromagnetic VIE for continuously varying media;
\item broadband electromagnetic internal states;
\item nonlinear magnetic materials;
\item internal conductor branching/junction networks;
\item phase change or radiative thermal nonlinearity;
\item learned operator-level preconditioner correction.
\end{enumerate}
The public Scene schema already has extension points for these features; they
are absent from the first backend, not removed from the theory.

\subsection{MVP physics pipeline}

The complete MVP path is
\[
 \text{continuous coils}
 \rightarrow
 \text{finite-cross-section current--charge IE}
 \rightarrow
 \{Z,H_q\}
 \rightarrow
 \text{loss channels}
 \rightarrow
 \text{stable thermal ROM}
 \rightarrow
 T(\mathbf r,t).
\]
The neural path is added only after the REFERENCE path is independently
converged.

\subsection{MVP reference benchmarks}

At minimum the teacher must pass the following locked tests.

\paragraph{DC resistance.}
For straight/sufficiently simple conductors in the low-frequency limit, compare
against
\[
 R_{\rm dc}=\rho_e L/A
\]
including convergence as modal rank is reduced toward the uniform mode.

\paragraph{Self inductance trend.}
For circular/superelliptic loops, compare against trusted analytic/asymptotic
limits where valid and an independent converged numerical oracle elsewhere.

\paragraph{Mutual inductance.}
For coaxial circular loops, compare against the classical elliptic-integral
reference. Then rotate/translate both loops together and require invariant port
results to numerical tolerance.

\paragraph{Reciprocity.}
For every reciprocal two/multi-coil scene,
\[
 Z-Z^\T
\]
must converge to zero under teacher refinement.

\paragraph{Passivity.}
For random complex currents,
\[
 \Rea(i^\ast Zi)\ge0
\]
up to declared numerical tolerance.

\paragraph{Skin/proximity convergence.}
Sweep frequency so
\[
 \min(w,h)/\delta
\]
crosses from $\ll1$ to $O(1)$ and larger. Conductor modal enrichment must
converge without changing any background discretization.

\paragraph{Near-separation stress.}
Bring two conductors close without contact. Near-singular quadrature must
converge and the result must remain independent of arbitrary quadrature
partition orientation.

\paragraph{Pose invariance.}
Apply random common rigid transforms to an entire scene. Global port matrices
must remain invariant; spatial current/loss queries must transform
equivariantly.

\paragraph{Power closure.}
For every port-current test,
\[
 \frac12 i^\ast\Herm(Z)i
\]
must agree with integrated conductor heat plus any declared outward channel.

\paragraph{Thermal impulse/step response.}
Compare Green/reduced thermal response against an independent high-accuracy
thermal reference for selected geometries and source channels.

\paragraph{Thermal energy balance.}
Over every accepted trajectory,
input heat minus environmental loss equals stored thermal-energy change to the
declared integration tolerance.

\subsection{MVP numerical acceptance targets}

Exact numerical values are configuration parameters, but the first development
target is:
\begin{center}
\begin{tabular}{ll}
\toprule
Quantity & Initial target\\
\midrule
teacher algebraic residual & $\le10^{-10}$\\
reciprocity defect & $\le10^{-10}$ relative\\
power closure & $\le10^{-6}$ relative\\
basis/quadrature convergence & $\le10^{-4}$ relative\\
independent port-$Z$ benchmark & $\le5\times10^{-3}$ relative\\
thermal transfer benchmark & $\le10^{-2}$ relative\\
FAST port-$Z$ validation & $\le10^{-2}$ median, $\le5\times10^{-2}$ max\\
FAST thermal observable validation & $\le5\times10^{-2}$ max\\
CERTIFIED requested observables & user tolerance or explicit uncertified status\\
\bottomrule
\end{tabular}
\end{center}
These are engineering development targets, not claims of guaranteed universal
accuracy. Release values may be tightened after the teacher is characterized.

\subsection{MVP neural model}

Only after the REFERENCE dataset passes the benchmarks above:
\begin{enumerate}
\item train isolated-object self models;
\item train pair interaction kernels over relative pose and intrinsic geometry;
\item form the analytic/self/pair baseline;
\item train a small many-body residual graph network;
\item train the continuous conductor-loss factor decoder;
\item train structured thermal coefficient corrections;
\item calibrate uncertainty against fixed validation and CERTIFIED residuals.
\end{enumerate}

For two-coil homogeneous scenes, the many-body residual may be nearly zero. In
that case the correct result is a small/simple neural model, not artificially
increasing network capacity.

\subsection{MVP FAST interface}

A minimal query is
\[
 \texttt{predict(scene, excitation, times, points, mode=FAST)}.
\]
It returns:
\begin{enumerate}
\item reciprocal/passive $Z(\omega)$;
\item currents/voltages after optional circuit solve;
\item total conductor loss and optional continuous $q(\mathbf r)$ queries;
\item temperatures at requested times/points;
\item uncertainty/calibration metadata.
\end{enumerate}

\subsection{MVP CERTIFIED interface}

CERTIFIED mode additionally returns:
\begin{enumerate}
\item electromagnetic residual and correction count;
\item port observable error estimate;
\item power-closure residual;
\item thermal state/integration residual;
\item final certification status.
\end{enumerate}

\subsection{MVP stop conditions}

Implementation stops and the theory is revisited if any of the following occur:
\begin{enumerate}
\item conductor self accuracy requires a world-volume mesh;
\item cross-section rank grows comparably to a full 3D conductor volume
      discretization over the target regime;
\item common rigid motion changes port results materially;
\item the teacher needs empirical fine-minus-coarse correction to pass basic
      conductor benchmarks;
\item FAST accuracy improves only by predicting world-coordinate fields;
\item CERTIFIED correction routinely costs as much as a cold REFERENCE solve.
\end{enumerate}

These are architecture falsification conditions, not bugs to patch around.

\subsection{MVP completion criterion}

The first implementation is complete only when the same artifact can:
\begin{enumerate}
\item evaluate previously unseen arbitrary-pose superelliptic coils;
\item produce passive reciprocal port response;
\item provide continuous conductor heat queries;
\item evolve temperature in time;
\item run FAST without a Maxwell solve;
\item escalate the same query to CERTIFIED mode;
\item reproduce locked REFERENCE benchmarks without geometry-specific hacks.
\end{enumerate}
Only then should package SIE/VIE and broadband physics become implementation
priorities.
